\begin{annexitem}{annex:p14.translation.standard-interface}{thm:p14-standard-translation-interface}{Truth and consequence of the standard translation}
\begin{proof}
The structure transformations use the same nonempty underlying set. Each
first-order binary or unary relation is exactly the corresponding modal
accessibility or atomic extension. Hence decoding an encoded Kripke structure
recovers those data, and encoding a decoded first-order structure recovers
its relations. No relation is chosen from a family of possible interpretations.

First show by induction that the only possible free object variable of
\(\operatorname{ST}_x(A)\) is \(x\). Atoms use only \(x\), and falsity
uses no variable. Negation and Boolean combinations preserve the property.
At the box clause, the recursively translated subformula at \(y\) has
only possible free variable \(y\); \(Rxy\) has free variables \(x,y\).
Binding the fresh \(y\) leaves only \(x\). The existential diamond
clause has the same free-variable calculation. Freshening the finite
formula's bound variables avoids every capture.

For truth preservation, proceed by induction on the modal formula,
simultaneously for every displayed free world variable and every total
assignment with its prescribed point value.
For an atom \(p\),
\[
 \mathcal M,w\entails p
 \Longleftrightarrow w\in V(p)
 \Longleftrightarrow s(x)\in P_p^{\mathfrak A_{\mathcal M}}
 \Longleftrightarrow\mathfrak A_{\mathcal M},s\entails P_p(x).
\]
Falsity is false in both semantics. Negation reverses the same inductively
equal truth value. For each of conjunction, disjunction and implication,
both semantics apply the same classical Boolean operation to the two
inductively equal operand values. This proves every Boolean clause, rather
than changing the source connective interpretation.

For the box, its modal truth is
\[
 \forall v\in W\ (wRv\Longrightarrow\mathcal M,v\entails A).
\]
For every such \(v\), the updated assignment \(s[y\mapsto v]\)
agrees with the pointed induction instance at the free variable \(y\).
The induction hypothesis makes this precisely
\[
 \forall v\in W\ (R(s(x),v)\Longrightarrow
       \mathfrak A_{\mathcal M},s[y\mapsto v]
                            \entails\operatorname{ST}_y(A)),
\]
which is the first-order universal clause. The variable \(y\) is fresh,
so updating it preserves the value of \(x\) in \(Rxy\). For diamond,
modal truth means that some \(v\) is an accessible world satisfying
\(A\). The induction hypothesis and the same update identify this
condition with the existential formula
\(\exists y(Rxy\land\operatorname{ST}_y(A))\).
This agrees as well with the classical dual of the box clause, because
diamond is the stipulated \(\neg\Box\neg\) abbreviation. The induction
proves the displayed satisfaction equivalence in every encoded model.
Decoding an arbitrary first-order structure gives the same equivalence
in the reverse direction. If two assignments agree at \(x\), both
translations therefore have the truth value of the same modal formula
at that same point, proving assignment independence directly.

For the consequence equivalence, suppose local modal consequence holds.
Let an arbitrary nonempty first-order structure and assignment satisfy
every member of \(\operatorname{ST}_x[\Gamma]\). Decode it to a Kripke
model pointed at \(s(x)\). The truth equivalence makes every modal premise
true there, so the modal consequence makes \(A\) true there, and the
same equivalence gives \(\operatorname{ST}_x(A)\). Conversely, suppose
the first-order consequence holds, and let a pointed Kripke model satisfy
every member of \(\Gamma\). Encode it and choose an assignment with
\(s(x)=w\). For example, sending every object variable to \(w\)
gives a total assignment without further choices. Every translated premise
holds, so the translated conclusion and then the modal conclusion hold.
The argument uses every member of an arbitrary \(\Gamma\) pointwise;
it does not require enumerating \(\Gamma\) or applying compactness.

If \(A\) is modally satisfiable, encode a satisfying pointed model. Its
world supplies a witness to \(\exists x\operatorname{ST}_x(A)\).
Conversely, a first-order structure satisfying that sentence supplies an
object witness; decode and point at it to satisfy \(A\).
For validity, if \(A\) holds at every point of every nonempty Kripke
model, then decoding any first-order structure makes the translation
true at each value of \(x\), so its universal closure is valid.
Conversely, encode any pointed model and apply the universally closed
valid formula at its point, then use the truth equivalence.
These establish both directions of both closure statements. Nonemptiness
is retained throughout, since a pointed modal model cannot use an empty
world set. None of the steps introduces a finite-model or Choice premise.
\end{proof}
\end{annexitem}

\begin{annexitem}{annex:p14.translation.selected-reachability}{thm:p14-selected-reachability}{Selected relation objects on a fixed finite graph}
\begin{proof}
The object domain has two distinct elements. Every selected relation
domain is nonempty. The predicate constant \(E\) denotes the empty
binary relation, which belongs to the displayed \(\mathcal R_2\).
Constants \(a,b\) have the stipulated distinct object values. Thus
\(\mathcal H\) satisfies every typing and selected-domain condition
of an arbitrary relational Henkin structure.

The exact reachability formula is
\[
 \forall X\bigl[(X(a)\land
     \forall u\forall v((X(u)\land E(u,v))\to X(v)))\to X(b)\bigr].
\]
In \(\mathcal H\), the only allowed unary \(X\) is \(M^1\).
It contains the one-entry tuples of both values of \(a,b\), and every closure implication
holds because \(E\) has no pairs. The displayed implication is
therefore true for the sole quantified relation, so Reach holds.

In \(\mathcal F\), every subset of \(M^1\) is available. Instantiate
\(X\) by \(\{(0)\}\). It contains the tuple of \(a\), is closed under the empty
edge relation, and does not contain the tuple of \(b\). The antecedent of the
implication is true and its consequent false. Thus the universal formula
fails. A positive-length directed path would require an edge, and a
length-zero path has equal endpoints. Neither exists from 0 to 1.

Finally take the eligible unary comprehension formula \(x=a\), whose
defining relation variable is fresh. The instance requires
\[
 \exists X\,\forall x\,(X(x)\leftrightarrow x=a).
\]
For the sole available \(X=M^1\) and an assignment \(x\mapsto1\),
\(X(x)\) is true while \(x=a\) is false, so the biconditional does not
hold at every object. No witness exists
in \(\mathcal R_1\), proving failure of that comprehension instance.
The matched class requires every eligible instance, so \(\mathcal H\)
does not belong to it. The object domain and each selected relation domain
are finite and explicitly specified; their family is defined uniformly
by arity. Neither a metatheoretic choice argument nor an identification
of full and selected relation ranges is used.
\end{proof}
\end{annexitem}
